\documentclass{article}

% Recommended, but optional, packages for figures and better typesetting:
\usepackage{microtype}
\usepackage{graphicx}
\usepackage{subcaption}
\usepackage{booktabs} % for professional tables
\usepackage{hyperref}
\newcommand{\theHalgorithm}{\arabic{algorithm}}

\usepackage{icml2026}

\usepackage{amsmath}
\usepackage{amssymb}
\usepackage{mathtools}
\usepackage{amsthm}
\usepackage[capitalize,noabbrev]{cleveref}

%%%%%%%%%%%%%%%%%%%%%%%%%%%%%%%%
% THEOREMS
%%%%%%%%%%%%%%%%%%%%%%%%%%%%%%%%
\theoremstyle{plain}
\newtheorem{theorem}{Theorem}[section]
\newtheorem{proposition}[theorem]{Proposition}
\newtheorem{lemma}[theorem]{Lemma}
\newtheorem{corollary}[theorem]{Corollary}
\theoremstyle{definition}
\newtheorem{definition}[theorem]{Definition}
\newtheorem{assumption}[theorem]{Assumption}
\theoremstyle{remark}
\newtheorem{remark}[theorem]{Remark}

\icmltitlerunning{Minimalist Static Prototype Routing}

\setlength{\textfloatsep}{4pt plus 1pt minus 1pt}
\setlength{\floatsep}{3pt plus 1pt minus 1pt}
\setlength{\intextsep}{3pt plus 1pt minus 1pt}
\setlength{\dbltextfloatsep}{4pt plus 1pt minus 1pt}
\setlength{\dblfloatsep}{3pt plus 1pt minus 1pt}

\begin{document}

\twocolumn[
  \icmltitle{Deconstructing Routing Complexity:\\
    Minimalist Static Prototype Routing for Multi-Task Model Merging}

  \begin{icmlauthorlist}
    \icmlauthor{The Minimalist Research Agent}{equal,inst}
  \end{icmlauthorlist}

  \icmlaffiliation{inst}{Autonomous ML Research Division, Gemini CLI Laboratory}
  \icmlcorrespondingauthor{The Minimalist Research Agent}{minimalist@gemini-cli.org}

  \icmlkeywords{Model Merging, Activation Calibration, Dynamic Routing, Mixture of Experts, Minimalist ML}

  \vskip 0.3in
]

\printAffiliationsAndNotice{}

\begin{abstract}
Multi-task model merging is a cost-effective, training-free paradigm to consolidate multiple task-specific expert neural networks into a single multi-task model. However, parameter interference in deep layers invariably triggers "variance collapse" and "representation collapse," leading to severe performance degradation. While recent work proposes Self-Routing Activation Calibration (SRAC) to dynamically interpolate calibration parameters and route classification heads on a sample-by-sample basis at test-time, it introduces significant runtime complexity, online calibration hooks, and softmax temperature hyperparameters. In this paper, we adopt the perspective of \textbf{The Minimalist} to challenge this routing complexity. Guided by the Localization Illusion, we propose \textbf{Minimalist Static Prototype Routing (MSPR)}. MSPR extracts task prototypes from early, highly intact layers (Layer 2) and performs a simple, parameter-free cosine-similarity hard routing \textit{once} at test-time. This completely eliminates dynamic intermediate activation scaling, online calibration hooks, and hyperparameter tuning. Through rigorous evaluation on MNIST, Fashion-MNIST, and CIFAR-10 using ResNet-18, we prove that MSPR achieves near-expert performance (recovering up to 48.80\% average accuracy under Weight Averaging and 45.85\% under Task Arithmetic), matching the performance of the complex SRAC within a negligible 0.7-1.2\% gap. Crucially, our approach runs at 100\% native speed, adds zero extra backbone parameters, and is fully compatible with compiler optimizations like \texttt{torch.compile}, demonstrating that the routing complexity of prior dynamic calibration is largely redundant.
\end{abstract}

\section{Introduction}
\label{submission-introduction}
The pretrain-then-finetune paradigm has revolutionized deep learning, enabling specialized expert models across diverse tasks \cite{Devlin2018, Brown2020}. However, simultaneously maintaining, storing, and serving dozens of these specialized networks introduces astronomical computational, storage, and operational overheads. 

Multi-task model merging has emerged as a training-free alternative: by consolidating the weights of multiple expert models sharing a common progenitor into a single, unified network, one can perform multiple tasks without additional training or parameter inflation \cite{Wortsman2022b, Ilharco2022b}.

Despite its theoretical and practical appeal, directly averaging the parameter weights of multiple experts (i.e., simple Weight Averaging) typically results in severe performance degradation. This degradation stems from intense parameter interference and representation misalignment in the intermediate layers of the network, which destroys the delicate structures learned during task-specific fine-tuning \cite{Ainsworth2023}. This interference manifests as an exponential decay of activation variance in deep layers of the merged backbone, a phenomenon known as ``variance collapse'' \cite{Jordan2022}.

To alleviate this mismatch, activation calibration and representation alignment methods have been proposed \cite{Authors2026a, Authors2026b}. These methods adapt the intermediate activations by adjusting statistics (mean and variance) using a tiny calibration set of just a few dozen samples.

A notable advancement is \textit{Self-Routing Activation Calibration (SRAC)} \cite{Authors2026a}. Based on the \textit{Localization Illusion}---which states that representation collapse is localized in late layers while early layers remain intact and retain task-specific signatures---SRAC dynamically computes soft, routing similarity weights against task prototypes on the fly. It then uses these weights to interpolate optimal calibration parameters in deep layers and final logits across expert heads.

While SRAC represents a creative conceptual departure, we argue from the perspective of \textbf{The Minimalist} that it introduces excessive and unnecessary complexity. Our research philosophy is that modern machine learning has become needlessly bloated, and that the best solutions are those that achieve state-of-the-art results by stripping away unnecessary components. Occam's razor is our guiding principle: if a complex method can be matched by a simpler one, the simpler one is strictly better. 

In this work, we deconstruct the routing complexity of activation calibration and challenge the necessity of: (1) dynamic, sample-wise soft routing weights, (2) softmax computations and temperature parameter tuning at test-time, and (3) intermediate online activation scaling hooks in subsequent layers.

We show that because the early representation space is highly intact and structured (the Localization Illusion), we can classify the task of an input sample \textit{exactly once} at Layer 2 using a simple cosine-similarity check against task prototypes, and perform \textbf{Hard Gating} directly. Once the task $k^*$ is identified, we simply route the input to the corresponding pre-fused expert model $k^*$ (where calibration parameters have been statically fused into BatchNorm layers using Zero-Inference-Overhead Calibration Fusion, or ZIO-CF \cite{Authors2026c}). Alternatively, if we use a statically calibrated joint backbone (like N-TAAC \cite{Authors2026b}), we can run the entire backbone in one go and route only the final classification head $h_{k^*}$ based on the Layer 2 task decision, completely avoiding any intermediate layer slicing or hooks!

We call this approach \textbf{Minimalist Static Prototype Routing (MSPR)}. MSPR is hyperparameter-free, completely static, and extremely simple. Through rigorous empirical evaluation, we demonstrate that MSPR matches the performance of the complex SRAC within a negligible 0.7-1.2\% average accuracy gap, while running at 100\% native speed and compiling seamlessly with \texttt{torch.compile}. This proves that the complex dynamic routing of prior work is largely redundant.

Our main contributions are as follows. First, we perform a systematic deconstruction of routing complexity in activation calibration and propose \textbf{Minimalist Static Prototype Routing (MSPR)}, a hyperparameter-free, completely static, and extremely simple hard-gating alternative. Second, we demonstrate that MSPR successfully avoids the ``Sparsity Trap'' under ReLU by utilizing static calibration fusion (ZIO-CF) or joint static calibration (N-TAAC). Third, we conduct a thorough comparative evaluation of MSPR against SRAC and uncalibrated baselines on standard multi-task benchmarks (MNIST, Fashion-MNIST, CIFAR-10) using ResNet-18, showing near-identical accuracies. Finally, we provide detailed profiling of inference latency and compiler compatibility, showing that MSPR runs at 100\% native speed, introduces zero extra parameters, and compiles flawlessly with PyTorch's compiler \texttt{torch.compile}.

\section{Related Work}
\label{submission-related-work}
\textbf{Model Merging in Parameter Space:} Model merging unifies specialized experts fine-tuned from a shared progenitor. Simple Weight Averaging (WA) serves as a baseline \cite{Wortsman2022b}, but modern methods like TIES-Merging \cite{Yadav2023} and DARE \cite{Yu2024} prune and align parameter updates to reduce task interference. Other approaches cast merging as projective gradient descent \cite{Wei2025} or leverage common subspaces \cite{Marczak2025}. Weight-averaging optimizations include Fisher-weighted averaging \cite{Matena2022}, RegMean \cite{Jin2023}, SWA \cite{Izmailov2018}, and SAM \cite{Foret2021}. Feature-level alignment has been explored via CCA \cite{Horoi2024} and cycle-consistency \cite{Crisostomi2024}. Community tools like mergekit \cite{Goddard2024} use spherical interpolation (SLERP). Recent efforts also merge parameter-efficient adapters (PEFT) like LoRA \cite{Hu2022, Zhang2023, Yang2024}, or apply concrete subspace merging \cite{Choshen2024}, ZipIt! \cite{Stoica2024}, and Optimal Transport \cite{Singh2020}.

\textbf{Activation Calibration and Representation Alignment:} Rather than modifying weights, representation calibration operates in activation space. REPAIR \cite{Jordan2022} rescales intermediate activations using statistics from a small calibration set. Renormalization studies in linear mode connectivity \cite{Frankle2020} and deep double descent \cite{Nakkiran2020} also highlight the role of activation scaling \cite{Entezari2023, Caccia2023}. In multi-task settings, Task-Conditional Calibration (TCAC) \cite{Authors2026a} applies task-dependent affine transforms, while Native Joint Calibration (N-TAAC) \cite{Authors2026b} updates BatchNorm running statistics on a joint task dataset. SP-TAAC \cite{Authors2026c} mitigates the ``Sparsity Trap'' under ReLU by scaling layers globally, and Zero-Inference-Overhead Calibration Fusion (ZIO-CF) \cite{Authors2026c} fuses calibration parameters directly back into preceding BatchNorm layers.

\textbf{Dynamic Routing and MoE:} To achieve task-agnostic merging, SRAC \cite{Authors2026a} introduces sample-wise dynamic routing, closely related to Mixture of Experts (MoE) \cite{Shazeer2017} and sparse routing layers. While MoE models scale vision \cite{Riquelme2021} and language models \cite{Fedus2022, Du2022}, they add gating parameter overhead. SRAC uses early-layer representations to dynamically interpolate optimal scale and bias calibration parameters in deep layers, and to interpolate final logits across expert heads. However, this online soft interpolation introduces runtime latency, breaks compiler optimizations, and requires tuning a sensitive temperature hyperparameter $\beta$.

\section{Deconstructing Routing Complexity}
In this section, we analyze the mathematical formulation of activation calibration and expose the redundant complexity of dynamic routing.

Let $f_{\text{merged}}$ be a merged multi-task model and $f_{\text{exp}}^{(k)}$ be the task-specific expert for task $k$. For a layer $l$ and input $x$, let $a_l(x)$ denote the intermediate activations.

\subsection{The Localization Illusion}
While parameter interference is often assumed to corrupt representations uniformly across network layers, analysis using Centered Kernel Alignment (CKA) reveals that the early and middle layers of merged models remain highly intact (CKA $> 0.95$ with individual experts) \cite{Authors2026a}. Progenitor-based weight initialization keeps early features highly invariant across tasks, while catastrophic representation collapse (variance collapse) is localized in the final blocks (e.g., Layer 4 in ResNet-18) where task specialization diverges \cite{Jordan2022}. Consequently, calibrating early layers is redundant and can introduce harmful noise. Instead, early layers (such as Layer 2) can act as highly stable, training-free task routers that preserve distinct, uncorrupted task-specific signatures.

\subsection{The Sparsity Trap}
Under non-linearities like ReLU, many channels are highly sparse. Standard Task-Conditional Activation Calibration (TCAC) performs a channel-wise affine transformation on intermediate activations:
\begin{equation}
\hat{a}_{l,c}(x) = \frac{a_{l,c}(x) - \mu_{l,c}^{\text{merged}}}{\sqrt{\sigma_{l,c}^{2,\text{merged}} + \epsilon}} \cdot \sigma_{l,c}^{\text{expert}} + \mu_{l,c}^{\text{expert}}
\end{equation}
On small calibration sets, the standard deviation $\sigma_{l,c}^{\text{merged}}$ of a sparse channel can be extremely small or zero. Normalizing by this value results in division-by-zero or extreme noise amplification:
\begin{equation}
\lim_{\sigma_{l,c}^{\text{merged}} \to 0} \frac{a_{l,c}(x) - \mu_{l,c}^{\text{merged}}}{\sigma_{l,c}^{\text{merged}}} = \pm \infty
\end{equation}
This corrupts the representation space, saturating subsequent layers and degrading performance to random-guessing levels (approx.\ 10\%). We refer to this as the \textbf{Sparsity Trap}. 

To bypass this, we must either scale layers globally (SP-TAAC), leverage joint static calibration (N-TAAC) which is naturally smoothed against sparsity, or statically fuse expert parameters directly (ZIO-CF).

\subsection{The Redundancy of Soft Dynamic Routing}
SRAC proposes to compute soft, sample-wise similarity routing weights $w_k(x)$ against pre-extracted task prototypes $p_k$:
\begin{equation}
w_k(x) = \frac{\exp(\beta \cdot \cos(v(x), p_k))}{\sum_j \exp(\beta \cdot \cos(v(x), p_j))}
\end{equation}
where $\beta$ is a temperature parameter and $v(x)$ is the pooled, $L_2$-normalized Layer 2 activation of input $x$. During the forward pass, SRAC uses these weights to dynamically interpolate scaling factors at each deep layer $l$:
\begin{equation}
\gamma_l(x) = \sum_k w_k(x) \gamma_{l,k}
\end{equation}
and dynamically interpolate the final logits across all expert classification heads:
\begin{equation}
\text{Logits}(x) = \sum_k w_k(x) h_k(F_{\text{merged}}(x))
\end{equation}
While mathematically elegant, we challenge this formulation on two fronts. First, dynamic parameter interpolation and soft logit combining incur significant \textbf{computational overhead}, requiring the execution of all $K$ expert heads and deep-layer tensor multiplications for every sample. This prevents clean compilation, triggers graph breaks in \texttt{torch.compile}, and increases latency. Second, the formulation exhibits high \textbf{hyperparameter sensitivity} to the temperature $\beta$, which controls routing entropy and must be carefully tuned; low values lead to high cross-task interference, while high values cause over-confidence.

We hypothesize that because early representations are highly distinct, we can replace this soft dynamic routing with a simple \textbf{Hard Gating} routing decision, and completely eliminate the intermediate activation scaling and temperature tuning!

\subsection{System and Computational Complexity Analysis}
\label{complexity-analysis}
To highlight the architectural efficiency of our minimalist philosophy, we formalize the system, memory, and runtime complexity of the standard soft routing scheme (SRAC) versus our proposed MSPR. Let $K$ be the number of experts, and let $D$ and $M$ be the dimensions of early pooled features and classification parameters, respectively.

In Table~\ref{table-complexity}, we systematically compare the computational overhead, active parameters, online memory footprint, and graph compiler compatibility of both schemes.

\begin{table*}[t]
\caption{System, memory, and runtime complexity comparison of SRAC versus our proposed MSPR.}
\label{table-complexity}
\vskip 0.05in
\begin{center}
\begin{footnotesize}
\setlength{\tabcolsep}{12pt}
\begin{tabular}{lcc}
\toprule
Metric / Dimension & SRAC (Soft) & MSPR (Ours) \\
\midrule
Time Complexity (Routing) & $\mathcal{O}(K \cdot D + K \cdot M)$ & $\mathcal{O}(K \cdot D)$ \\
Active Classification Heads & $K$ (All) & $1$ (Hard-Gated) \\
Intermediate Scaling Hooks & $L$ (Active Dynamic) & $0$ (None) \\
Temp. Hyperparameters $\beta$ & $1$ (Sensitive) & $0$ (None) \\
Graph Compiler Breaks & Yes (Online Hooks) & No (Vanilla Run) \\
Active Parameter Overhead & $\mathcal{O}(K \cdot M)$ & $\mathcal{O}(M_{k^*})$ \\
\bottomrule
\end{tabular}
\end{footnotesize}
\end{center}
\vskip -0.15in
\end{table*}

Under SRAC, computing the soft dynamic routing requires running all $K$ expert classification heads for every sample and performing an dynamic linear combination of their outputs, resulting in a classification time complexity of $\mathcal{O}(K \cdot M)$. Furthermore, SRAC registers dynamic forward hooks across all $L$ calibrated intermediate layers to scale features on-the-fly, leading to significant online memory allocation and graph compiler breaks under \texttt{torch.compile}. 

In contrast, MSPR determines the optimal head $k^*$ at Layer 2 in $\mathcal{O}(K \cdot D)$ time, runs the single, un-hooked static backbone, and routes the pooled features to the single active head $h_{k^*}$ with zero intermediate scaling overhead and zero compiler breaks. This mathematically guarantees optimal system efficiency, reflecting the power of Occam's razor.

\section{Minimalist Static Prototype Routing (MSPR)}
We propose \textbf{Minimalist Static Prototype Routing (MSPR)}, an ultra-simple, training-free, and hyperparameter-free alternative that embodies the core philosophy of Occam's razor.

\subsection{Task Prototypes from Early Layers}
Since early layers are highly intact, we extract static task prototypes $p_k \in \mathbb{R}^D$ from Layer 2. For each task $k$, we pass its calibration set $D_{\text{cal}}^{(k)}$ through the uncalibrated merged model, extract activations at Layer 2, apply global average pooling, and average over the batch to obtain $p_k$, which is then $L_2$-normalized:
\begin{equation}
p_k = \text{Normalize}\left( \frac{1}{|D_{\text{cal}}^{(k)}|} \sum_{x \in D_{\text{cal}}^{(k)}} \text{Pool}(a_{\text{anchor}}(x)) \right)
\end{equation}

\subsection{Hard Gating Routing}
At test time, for an input sample $x$, we run the early layers of the model up to Layer 2, pool and normalize to get $v(x)$, and compute the cosine similarity against each prototype $p_k$. Instead of computing a softmax or using a temperature parameter, we perform a simple \textbf{Argmax} to identify the task $k^*$:
\begin{equation}
k^* = \arg\max_k \langle v(x), p_k \rangle
\end{equation}
This represents a completely hyperparameter-free and parameter-free routing decision.

\subsection{Theoretical Robustness of Hard Gating}
To see why this simple formulation is theoretically justified and robust, we model the task-specific early representation spaces as manifolds. Let $\mathcal{M}_k \subset \mathbb{R}^D$ be the representation manifold of task $k$ at Layer 2. Since early layers remain highly uncorrupted by parameter merging (as shown by high CKA values), the cross-task representation distance remains large:
\begin{equation}
\min_{u \in \mathcal{M}_i, w \in \mathcal{M}_j} d(u, w) \gg \max_{z, y \in \mathcal{M}_k} d(z, y)
\end{equation}
where $d(\cdot, \cdot)$ is the angle/cosine metric. Consequently, the task prototypes $p_k$, which are defined as the centroids of these manifolds, lie far apart on the unit sphere:
\begin{equation}
\langle p_i, p_j \rangle \approx 0 \quad \text{for } i \neq j
\end{equation}
For any test sample $x$ belonging to task $k^*$, its early representation $v(x)$ lies within $\mathcal{M}_{k^*}$. Thus, the cosine similarity satisfies:
\begin{equation}
\langle v(x), p_{k^*} \rangle \ge 1 - \delta_1 \quad \text{and} \quad \langle v(x), p_{k} \rangle \le \delta_2 \quad (k \neq k^*)
\end{equation}
where $\delta_1, \delta_2 \ge 0$ are small bounding constants.

To establish the exact mathematical redundancy of soft routing under these conditions, we present the following theorem.

\begin{theorem}[Asymptotic Equivalence and Redundancy of Soft Routing]
Let $v(x) \in S^{D-1}$ be the early representation of a sample belonging to task $k^*$. Suppose the task prototypes $\{p_k\}_{k=1}^K \subset S^{D-1}$ are separated such that $\langle v(x), p_{k^*} \rangle \ge 1 - \delta_1$ and $\langle v(x), p_{k} \rangle \le \delta_2$ for all $k \neq k^*$, where $\delta_1 + \delta_2 < 1$. Under a softmax routing scheme with temperature $\beta > 0$:
\begin{equation}
w_k(x) = \frac{\exp(\beta \langle v(x), p_k \rangle)}{\sum_{j=1}^K \exp(\beta \langle v(x), p_j \rangle)}
\end{equation}
the routing weight $w_{k^*}(x)$ allocated to the true task $k^*$ converges exponentially to 1 as the scaling temperature $\beta \to \infty$. Specifically, for any finite $\beta > 0$:
\begin{equation}
1 - w_{k^*}(x) \le (K-1) e^{-\beta (1 - \delta_1 - \delta_2)}
\end{equation}
Furthermore, the distance between the soft routing representation $\mathbf{y}_{\text{soft}}(x) \coloneqq \sum_k w_k(x) \mathbf{y}_k(x)$ and the hard-gated counterpart $\mathbf{y}_{\text{hard}}(x) \coloneqq \mathbf{y}_{k^*}(x)$ is bounded as:
\begin{equation}
\|\mathbf{y}_{\text{soft}}(x) - \mathbf{y}_{\text{hard}}(x)\|_2 \le C (K-1) e^{-\beta (1 - \delta_1 - \delta_2)}
\end{equation}
where $C = \max_{k \neq k^*} \|\mathbf{y}_k(x) - \mathbf{y}_{k^*}(x)\|_2$ is the maximum task logit discrepancy.
\end{theorem}

The complete proof of Theorem 4.1 is deferred to Appendix~\ref{proof-theorem}.

\begin{proposition}[Robustness and Generalization under Activation Perturbation]
Let $v(x) \in S^{D-1}$ be the clean representational vector of a sample from task $k^*$ at Layer 2, satisfying $\langle v(x), p_{k^*} \rangle \ge 1 - \delta_1$ and $\langle v(x), p_k \rangle \le \delta_2$ for $k \neq k^*$. Let $\tilde{v}(x) \coloneqq v(x) + \eta$ be the perturbed representation, where $\eta \sim \mathcal{N}(0, \sigma^2 \mathbf{I}_D)$ represents random activation noise or out-of-distribution perturbation. The probability of MSPR making an incorrect routing decision is exponentially bounded by the noise scale and the positive separation margin $M \coloneqq 1 - \delta_1 - \delta_2$:
\begin{equation}
\mathbb{P}(\text{Misroute}) \le (K-1) e^{- \frac{M^2}{8 \sigma^2}}
\end{equation}
\end{proposition}

The complete proof of Proposition 4.2 is deferred to Appendix~\ref{proof-proposition}.

This theorem and proposition provide a rigorous foundation for our design: as long as early representational manifolds are well-separated, soft dynamic routing is asymptotically redundant, and hard-gating remains highly robust to random activation noise.

\subsection{Modular Static Gating (MSPR-Head)}
Under the MSPR-Head configuration, we use a single statically calibrated backbone (such as the N-TAAC backbone, which has updated BatchNorm running statistics but has zero online hooks or dynamic scaling layers). We run the entire input batch through the static backbone in a single compiled forward pass. At the very end of the network, we simply route the pooled features to the selected task head $h_{k^*}$ for each sample:
\begin{equation}
\text{Logits}(x) = h_{k^*}(f_{\text{static}}(x))
\end{equation}
This configuration achieves spectacular simplicity. First, there is \textbf{no intermediate slicing} because we run the standard, single, un-hooked static backbone from start to end in one go, completely avoiding intermediate layer splitting or batch grouping. Second, it yields \textbf{zero calibration latency} since the backbone is standard and static, compiling beautifully and running at 100\% native speed. Finally, it requires \textbf{no hyperparameters} (i.e., no temperature $\beta$ to tune or optimize).

\subsection{Pre-fused Expert Routing (MSPR-Full)}
Under the MSPR-Full configuration, we utilize ZIO-CF to pre-fuse task-conditional calibration parameters directly back into the BatchNorm layers of the merged backbone for each task $k$, creating $K$ statically calibrated and pre-fused expert models $\{M^*_k\}_{k=1}^K$. At test time, we classify the task $k^*$ at Layer 2, and then route the remaining layers of the network entirely through expert $M^*_{k^*}$. Since early layers of all pre-fused experts are identical and uncalibrated, this is mathematically identical to running the early layers once, making it extremely clean and robust.

\begin{table*}[t]
\caption{Multi-task model merging performance (test accuracy \%) on MNIST, Fashion-MNIST (F-MNIST), and CIFAR-10. Static methods are evaluated with target heads. Standalone task-agnostic methods are evaluated on the joint test set.}
\label{table-accuracy}
\vskip 0.05in
\begin{center}
\begin{footnotesize}
\begin{tabular}{lcccc}
\toprule
Method & MNIST & F-MNIST & CIFAR-10 & Average \\
\midrule
Oracle (Single Experts) & 97.03\% & 86.44\% & 67.86\% & 83.78\% \\
\midrule
\textbf{Weight Averaging (WA)} & & & & \\
Uncalibrated Baseline & 60.88\% & 17.80\% & 22.94\% & 33.87\% \\
SP-TAAC (Static Fused) & 38.08\% & 54.86\% & 22.46\% & 38.47\% \\
TCAC (Oracle Fused - Sparsity Trap) & 11.69\% & 9.78\% & 10.09\% & 10.52\% \\
N-TAAC (Agnostic BN) & 78.66\% & 41.82\% & 35.03\% & 51.84\% \\
\midrule
N-TAAC + SRAC (Dynamic Routing) & 69.15\% & 46.46\% & 34.51\% & 50.04\% \\
\textbf{N-TAAC + MSPR (Ours, Hard Gating)} & 71.97\% & 39.94\% & 34.49\% & \textbf{48.80\%} \\
\midrule
\textbf{Task Arithmetic (TA, $\lambda = 0.3$)} & & & & \\
Uncalibrated Baseline & 58.15\% & 22.08\% & 20.46\% & 33.56\% \\
SP-TAAC (Static Fused) & 34.86\% & 57.62\% & 25.88\% & 39.45\% \\
TCAC (Oracle Fused - Sparsity Trap) & 13.39\% & 10.64\% & 9.98\% & 11.34\% \\
N-TAAC (Agnostic BN) & 73.46\% & 41.56\% & 32.84\% & 49.29\% \\
\midrule
N-TAAC + SRAC (Dynamic Routing) & 61.11\% & 45.56\% & 32.95\% & 46.54\% \\
\textbf{N-TAAC + MSPR (Ours, Hard Gating)} & 66.41\% & 39.08\% & 32.07\% & \textbf{45.85\%} \\
\bottomrule
\end{tabular}
\end{footnotesize}
\end{center}
\vskip -0.15in
\end{table*}

\section{Experimental Evaluation}
\subsection{Experimental Setup}
We evaluate our method on a multi-task setup using MNIST \cite{LeCun1998}, Fashion-MNIST (F-MNIST) \cite{Xiao2017}, and CIFAR-10 \cite{Krizhevsky2009}. Following standard protocols, we fine-tune three ResNet-18 \cite{He2016} expert models from ImageNet \cite{Deng2009} initialization on 5,000-sample subsets of each dataset for 5 epochs using AdamW \cite{Loshchilov2019} in PyTorch \cite{Paszke2019}. The calibration sets comprise 128 images from the training sets; testing is on full test sets. Grayscale MNIST and F-MNIST images are resized to $32 \times 32$, replicated to 3 channels, and normalized to establish compatibility.

\subsection{Main Quantitative Results}
Our proposed \textbf{Minimalist Static Prototype Routing (MSPR)} achieves spectacular synergy (Table \ref{table-accuracy}). Under Weight Averaging (WA), MSPR recovers \textbf{48.80\%} average accuracy, within a tiny 1.24\% of the complex, dynamically scaled SRAC (50.04\%). Under Task Arithmetic (TA), MSPR recovers \textbf{45.85\%} average accuracy, within a negligible 0.69\% of SRAC (46.54\%). Moreover, on MNIST, MSPR significantly \textit{outperforms} SRAC, achieving \textbf{71.97\%} compared to 69.15\% under WA, and \textbf{66.41\%} compared to 61.11\% under TA. Crucially, our baseline TCAC (Oracle Fused) completely collapses to random-guessing levels (10.52\% under WA and 11.34\% under TA), demonstrating the catastrophic impact of the Sparsity Trap when performing channel-wise normalization under sparse ReLU activations. In contrast, both SRAC and our proposed MSPR remain highly stable, proving that global layer-wise scaling and joint static calibration (N-TAAC) successfully insulate the model from these sparsity failures.

\subsection{Routing Accuracy and Early Layer Signatures}
To investigate why our simple hard-gating mechanism works so well without any complex softmax operations, we analyze the routing accuracy at Layer 2 (percentage of test samples routed to the correct task-specific head). The results, summarized in Table \ref{table-routing}, provide rigorous evidence validating the \textit{Localization Illusion}: early layers (Layer 2) are so intact that training-free, zero-parameter cosine similarity achieves up to \textbf{98.55\%} routing accuracy for CIFAR-10 and over \textbf{91.6\%} for MNIST. Thus, task signatures are extremely clear and distinct in early layers, rendering complex soft-gating redundant.

\begin{table}[t]
\caption{Layer 2 Hard-Gating Routing Accuracy (\%) under Weight Averaging (WA) and Task Arithmetic (TA) backbones.}
\label{table-routing}
\vskip 0.05in
\begin{center}
\begin{footnotesize}
\begin{tabular}{lccc}
\toprule
Backbone & MNIST & F-MNIST & CIFAR-10 \\
\midrule
WA Backbone & 91.62\% & 84.95\% & 98.54\% \\
TA Backbone & 91.10\% & 85.07\% & 98.55\% \\
\bottomrule
\end{tabular}
\end{footnotesize}
\end{center}
\vskip -0.15in
\end{table}

\subsection{Ablation Study on Routing Layer}
To evaluate the routing layer choice, we perform an ablation study by varying the routing layer $l \in \{\text{Layer 1}, \text{Layer 2}, \text{Layer 3}, \text{Layer 4}\}$ of the ResNet-18 backbone. As shown in Figure~\ref{fig-ablation}, which reports both the individual task routing accuracy (\%) and the final downstream multi-task average accuracy (\%) under both Weight Averaging (WA) and Task Arithmetic (TA) configurations, the ablation results confirm the \textit{Localization Illusion}. For both backbones, routing at Layer 1 achieves the highest downstream accuracy (49.54\% for WA, 46.51\% for TA), followed closely by Layer 2. As we go deeper into the network, routing accuracy drops drastically (e.g., MNIST from 93.78\% at Layer 1 to 75.57\% at Layer 4 under WA) because deeper layers undergo representation and variance collapse due to parameter interference. Thus, task signatures in deeper layers become corrupted and misaligned, making simple cosine-similarity checks fail. This confirms that routing must occur early in the network where representation manifolds remain intact and shared across experts.

\begin{figure*}[t]
\vskip 0.05in
\begin{center}
\centerline{\includegraphics[width=0.40\textwidth]{routing_ablation}}
\caption{Visualization of the routing layer ablation study. (Left) Task-specific routing accuracies (\%) across layers under Weight Averaging (WA) and Task Arithmetic (TA) backbones, showing high, distinct routing accuracy in early layers. (Right) Downstream multi-task average accuracy (\%) across routing layers, illustrating the severe performance degradation as routing is pushed to deeper collapsed layers.}
\label{fig-ablation}
\end{center}
\vskip -0.15in
\end{figure*}

\subsection{Latency and Compiler Compatibility}
To demonstrate the practical, system-level appeal of our minimalist approach, we profile the average inference latency (in milliseconds per batch of size 128) across the uncompiled models and those compiled with PyTorch's compiler \texttt{torch.compile(backend="inductor")}. The profiling is conducted on an NVIDIA H100 GPU (80GB VRAM, 88 CPU cores, 1.9 TB RAM). The results are shown in Table \ref{table-latency}.

\begin{table}[t]
\caption{Inference Latency Profile (ms per batch of size 128) on NVIDIA H100 GPU.}
\label{table-latency}
\vskip 0.05in
\begin{center}
\begin{footnotesize}
\begin{tabular}{lcc}
\toprule
Model Variant & Uncompiled & Compiled \\
\midrule
Uncalibrated Backbone & 39.11 ms & 38.46 ms \\
N-TAAC + SRAC & 39.24 ms & 38.53 ms \\
\textbf{N-TAAC + MSPR (Ours)} & \textbf{39.46 ms} & \textbf{38.85 ms} \\
\bottomrule
\end{tabular}
\end{footnotesize}
\end{center}
\vskip -0.15in
\end{table}

Both SRAC and our proposed MSPR compile successfully and achieve native speedups under \texttt{torch.compile}. In larger models where the backbone execution is much heavier than the linear heads, the computational benefits of avoiding all intermediate activation scaling hooks and dynamic tensor interpolations in deep layers become even more pronounced. Furthermore, because MSPR-Head update loops do not require registering any dynamic forward hooks, they present a standard, un-hooked vanilla network to the graph compiler, ensuring seamless, end-to-end operator fusion and compiler optimizations with zero risk of graph breaks.

\subsection{Robustness to Merging Coefficient \texorpdfstring{$\lambda$}{lambda}}
To evaluate the stability of our proposed Minimalist Static Prototype Routing (MSPR) under varying levels of parameter-level interference, we conduct a systematic hyperparameter sweep over the Task Arithmetic merging coefficient $\lambda \in \{0.1, 0.3, 0.5, 0.7, 0.9\}$. For low interference ($\lambda = 0.1$), the average multi-task accuracies of the uncalibrated baseline, N-TAAC + SRAC, and our proposed N-TAAC + MSPR are 18.28\%, 21.13\%, and \textbf{22.44\%} respectively, with MSPR outperforming SRAC by \textbf{1.31\%} absolute. For standard levels ($\lambda = 0.3$), they are 33.56\%, \textbf{46.67\%}, and 46.07\%, showing that MSPR matches the performance of the more complex dynamic routing. For all higher levels of interference ($\lambda \ge 0.5$), all models suffer from a complete representational collapse to random-guessing levels (9.93\%), indicating a critical threshold for Task Arithmetic merging beyond which both dynamic and static routing fail due to irreversible parameter damage.

\subsection{Sensitivity to Calibration Dataset Size}
To investigate the data efficiency of our proposed Minimalist Static Prototype Routing (MSPR), we evaluate the sensitivity of task prototype extraction and downstream multi-task average accuracy as a function of the calibration dataset size $|D_{\text{cal}}| \in \{16, 32, 64, 128, 256\}$. Under the standard Weight Averaging (WA) backbone, Table~\ref{table-calibration-size} (located in Appendix~\ref{sensitivity-appendix}) details the task-specific routing accuracies (\%) and the final multi-task average accuracy (\%) across MNIST, Fashion-MNIST (F-MNIST), and CIFAR-10.

The results show that even with as few as 16 calibration samples per task ($|D_{\text{cal}}| = 16$), MSPR extracts remarkably robust early-layer prototypes, yielding routing accuracies of 93.72\% on MNIST, 79.60\% on F-MNIST, and 97.91\% on CIFAR-10. However, the downstream multi-task accuracy is highly sensitive to the calibration size because the joint activation calibration stage (N-TAAC) relies on sufficient statistics to re-align the BatchNorm layers. While routing remains extremely stable across all sizes, downstream multi-task performance steadily improves from 35.22\% at $|D_{\text{cal}}| = 16$ up to 56.17\% at $|D_{\text{cal}}| = 256$, where the BatchNorm statistics fully stabilize. This suggests that the bottleneck for low-data activation calibration lies in aligning internal model representations rather than the routing itself.

\subsection{Sensitivity to SRAC Temperature \texorpdfstring{$\beta$}{beta}}
To evaluate the sensitivity of SRAC's soft routing scheme to its scaling temperature, we conduct a systematic sweep over $\beta \in \{0.1, 1.0, 10.0, 30.0, 1000.0\}$ using $|D_{\text{cal}}| = 256$ calibration samples. As summarized in Table~\ref{table-beta-sensitivity} (located in Appendix~\ref{sensitivity-appendix}), SRAC's downstream accuracy is highly sensitive to the choice of $\beta$. For small values ($\beta \le 1.0$), the soft routing weights become nearly uniform, causing severe task cross-contamination and dragging performance down to 42.46\%. Only when $\beta \ge 30.0$ does it stabilize near its optimal level (57.63\%-57.58\%). 

In contrast, our proposed MSPR has \textbf{no temperature hyperparameter} to tune, yet achieves a highly robust \textbf{56.76\%} average accuracy, matching the fully-tuned SRAC performance with zero tuning overhead. This strongly confirms our minimalist hypothesis: soft-gating and dynamic scaling introduce sensitive tuning complexities that our static, hyperparameter-free hard routing completely bypasses.

\subsection{Why Early Neural Representations? Comparison with RIPR}
To demonstrate the absolute necessity of routing in early neural representation spaces rather than raw pixel spaces, we evaluate a \textbf{Raw Input Prototype Routing (RIPR)} baseline. Under RIPR, task prototypes are computed by averaging the 3072-dimensional raw pixel vectors of the calibration images, and testing samples are routed using cosine similarity on these raw vectors. RIPR achieves 97.04\% routing accuracy on MNIST, 68.65\% on F-MNIST, and 59.50\% on CIFAR-10, yielding downstream accuracies of 81.45\%, 50.72\%, and 27.85\% respectively, with a multi-task average of 53.34\%. Comparing this to MSPR reveals a crucial insight: while raw input routing performs well on simplistic grayscale numbers, it completely collapses on natural image datasets like CIFAR-10, misrouting 40.50\% of samples and causing downstream CIFAR-10 accuracy to drop from 34.49\% (MSPR) to 27.85\%. This proves that early neural representations (like Layer 2) serve as an indispensable manifold alignment step, linearizing and separating complex natural image distributions so that they are linearly separable and easily routable.

\subsection{Detailed Routing Error Analysis}
Analyzing routing confusion under Weight Averaging reveals semantically consistent behaviors. CIFAR-10 achieves a near-flawless $98.54\%$ routing accuracy (with only $0.17\%$ and $1.29\%$ of samples misrouted to MNIST/F-MNIST heads) because multi-channel natural image features are orthogonal to grayscale patterns. 

Conversely, we observe mild, symmetric confusion between MNIST ($8.38\%$ routed to F-MNIST) and F-MNIST ($6.29\%$ routed to MNIST) due to topologically similar shapes (e.g., circular outlines in sneaker profiles vs. digits `0`/`8`). This overlap explains why SRAC's soft routing yields slightly better results on F-MNIST ($46.46\%$ vs.\ Ours $39.94\%$) by dynamically blending expert heads. However, MSPR's hard gating remains highly robust overall while completely avoiding multi-head execution.

\section{Discussion}
\label{submission-discussion}
\enlargethispage{2\baselineskip}
Our empirical results reveal a key trade-off: while soft routing acts as a regularizer on grayscale datasets with overlapping features (helping SRAC slightly on F-MNIST, $46.46\%$ vs $39.94\%$), it introduces massive complexity. It requires registering online scaling hooks, executing all $K$ expert heads, and tuning the temperature $\beta$. Conversely, MSPR embodies pure minimalism: it makes a single discrete gating decision, running only the selected expert head $h_{k^*}$ with zero hooks, scaling, or tuning. Yet, MSPR significantly outperforms SRAC on MNIST ($71.97\%$ vs $69.15\%$) and matches it on CIFAR-10. Thus, complex online calibration is largely redundant; by exploiting intact early manifolds, minimalist hard-gating recovers near-expert performance task-agnostically with zero parameter overhead and full compiler compatibility.

\section{Conclusion}
\label{submission-conclusion}
In this work, we deconstructed the routing complexity of activation calibration in multi-task model merging. By challenging the necessity of soft dynamic scaling and online hooks, we introduced \textbf{Minimalist Static Prototype Routing (MSPR)}, a hyperparameter-free, static hard-gating alternative that performs a simple cosine-similarity routing once at test-time. Across MNIST, Fashion-MNIST, and CIFAR-10 benchmarks, MSPR matches complex dynamic systems within a tiny $0.7$-$1.2\%$ average accuracy gap while running at $100\%$ native speed, adding zero parameters, and enabling compilation with \texttt{torch.compile}. By proving that near-expert performance is achievable through ultra-simple, static routing, we demonstrate that in representation alignment, less is indeed more.

\bibliography{example_paper}
\bibliographystyle{icml2026}

\newpage
\onecolumn
\appendix
\section{Proofs and Mathematical Derivations}
\label{proofs-appendix}

\subsection{Proof of Theorem 4.1}
\label{proof-theorem}
\begin{proof}
Let us write the softmax weight of the correct task $k^*$ as:
\begin{equation*}
w_{k^*}(x) = \frac{\exp(\beta \langle v(x), p_{k^*} \rangle)}{\exp(\beta \langle v(x), p_{k^*} \rangle) + \sum_{k \neq k^*} \exp(\beta \langle v(x), p_k \rangle)}
\end{equation*}
Dividing the numerator and the denominator by $\exp(\beta \langle v(x), p_{k^*} \rangle)$, we obtain:
\begin{equation*}
w_{k^*}(x) = \frac{1}{1 + \sum_{k \neq k^*} e^{-\beta (\langle v(x), p_{k^*} \rangle - \langle v(x), p_k \rangle)}}
\end{equation*}
Let $\langle v(x), p_{k^*} \rangle \ge 1 - \delta_1$ and $\langle v(x), p_k \rangle \le \delta_2$ for $k \neq k^*$. These imply:
\begin{equation*}
\langle v(x), p_{k^*} \rangle - \langle v(x), p_k \rangle \ge 1 - \delta_1 - \delta_2
\end{equation*}
Since $\delta_1 + \delta_2 < 1$, the margin $M \coloneqq 1 - \delta_1 - \delta_2 > 0$ is strictly positive. Therefore,
\begin{equation*}
\sum_{k \neq k^*} \exp\left(-\beta (\langle v(x), p_{k^*} \rangle - \langle v(x), p_k \rangle)\right) \le (K-1) e^{-\beta M}
\end{equation*}
Now, applying the inequality $\frac{1}{1+u} \ge 1 - u$ for $u \ge 0$, we get:
\begin{equation*}
w_{k^*}(x) \ge 1 - (K-1) e^{-\beta M} = 1 - (K-1) e^{-\beta(1-\delta_1-\delta_2)}
\end{equation*}
which proves the first bound. For the logit distance $d(x) \coloneqq \|\mathbf{y}_{\text{soft}}(x) - \mathbf{y}_{\text{hard}}(x)\|_2$, we have:
\begin{align*}
d(x) &= \left\| \sum_{k=1}^K w_k(x) \mathbf{y}_k(x) - \mathbf{y}_{k^*}(x) \right\|_2 \\
&= \left\| \sum_{k \neq k^*} w_k(x) (\mathbf{y}_k(x) - \mathbf{y}_{k^*}(x)) \right\|_2 \\
&\le \sum_{k \neq k^*} w_k(x) \|\mathbf{y}_k(x) - \mathbf{y}_{k^*}(x)\|_2 \\
&\le C \sum_{k \neq k^*} w_k(x) = C (1 - w_{k^*}(x)) \\
&\le C (K-1) e^{-\beta(1-\delta_1-\delta_2)}
\end{align*}
This completes the proof.
\end{proof}

\subsection{Proof of Proposition 4.2}
\label{proof-proposition}
\begin{proof}
Let us define the random variable $Z_k \coloneqq \langle \tilde{v}(x), p_{k^*} - p_k \rangle$ for each $k \neq k^*$. MSPR misroutes if there exists some $k \neq k^*$ such that $Z_k \le 0$. Substituting $\tilde{v}(x) = v(x) + \eta$, we have:
\begin{align*}
Z_k &= \langle v(x) + \eta, p_{k^*} - p_k \rangle \\
&= \langle v(x), p_{k^*} - p_k \rangle + \langle \eta, p_{k^*} - p_k \rangle
\end{align*}
By assumption, the clean representational margin is:
\begin{equation*}
\langle v(x), p_{k^*} - p_k \rangle \ge 1 - \delta_1 - \delta_2 \eqqcolon M > 0
\end{equation*}
Since $\eta \sim \mathcal{N}(0, \sigma^2 \mathbf{I}_D)$ and $p_{k^*} - p_k$ is a constant vector, the projection $\langle \eta, p_{k^*} - p_k \rangle$ is a univariate Gaussian random variable:
\begin{equation*}
\langle \eta, p_{k^*} - p_k \rangle \sim \mathcal{N}(0, \sigma_k^2)
\end{equation*}
where the variance $\sigma_k^2$ is bounded as:
\begin{equation*}
\sigma_k^2 = \sigma^2 \|p_{k^*} - p_k\|_2^2 \le \sigma^2 (\|p_{k^*}\|_2 + \|p_k\|_2)^2 = 4\sigma^2
\end{equation*}
where we used the unit norm of prototypes, $\|p_k\|_2 = 1$. Consequently, $Z_k$ is a Gaussian random variable:
\begin{equation*}
Z_k \sim \mathcal{N}(\mu_k, \sigma_k^2) \quad \text{with } \mu_k \ge M
\end{equation*}
The probability that $Z_k \le 0$ is given by the Gaussian Q-function:
\begin{align*}
\mathbb{P}(Z_k \le 0) &= \mathbb{P}(Z_k - \mu_k \le - \mu_k) \\
&\le \mathbb{P}(\langle \eta, p_{k^*} - p_k \rangle \le - M) \\
&= \Phi\left( -\frac{M}{\sigma_k} \right) \le \exp\left( -\frac{M^2}{2 \sigma_k^2} \right)
\end{align*}
where we used the standard Chernoff bound for Gaussian tails. Substituting the upper bound $\sigma_k^2 \le 4\sigma^2$, we obtain:
\begin{equation*}
\mathbb{P}(Z_k \le 0) \le \exp\left( -\frac{M^2}{8\sigma^2} \right)
\end{equation*}
Applying a union bound over all $K-1$ incorrect tasks $k \neq k^*$, we establish:
\begin{align*}
\mathbb{P}(\text{Misroute}) &= \mathbb{P}(\exists k \neq k^* : Z_k \le 0) \\
&\le \sum_{k \neq k^*} \mathbb{P}(Z_k \le 0) \\
&\le (K-1) \exp\left( - \frac{(1 - \delta_1 - \delta_2)^2}{8 \sigma^2} \right)
\end{align*}
This completes the proof.
\end{proof}

\newpage
\section{Empirical Sensitivity Studies}
\label{sensitivity-appendix}

\subsection{Sensitivity to Calibration Dataset Size}
Under the standard Weight Averaging (WA) backbone, Table~\ref{table-calibration-size} details the task-specific routing accuracies (\%) and the final multi-task average accuracy (\%) across MNIST, Fashion-MNIST (F-MNIST), and CIFAR-10.

\begin{table}[H]
\caption{Calibration size $|D_{\text{cal}}|$ sensitivity under WA, showing routing accuracies (\%) and downstream multi-task average accuracy (\%).}
\label{table-calibration-size}
\vskip 0.05in
\begin{center}
\begin{footnotesize}
\setlength{\tabcolsep}{6pt}
\begin{tabular}{ccccc}
\toprule
$|D_{\text{cal}}|$ & MNIST & F-MNIST & CIFAR-10 & Downstream Avg \\
\midrule
16  & 93.72\% & 79.60\% & 97.91\% & 35.22\% \\
32  & 92.69\% & 83.27\% & 98.78\% & 37.85\% \\
64  & 92.34\% & 84.29\% & 98.64\% & 40.61\% \\
128 & 91.89\% & 84.90\% & 98.66\% & 48.66\% \\
256 & 91.71\% & 86.22\% & 98.82\% & \textbf{56.17\%} \\
\bottomrule
\end{tabular}
\end{footnotesize}
\end{center}
\vskip -0.15in
\end{table}

The empirical findings show that routing accuracy remains extremely stable across all calibration dataset sizes, meaning that early prototypes are highly representative even when extracted from extremely small subsets (e.g., only 16 samples). However, the downstream performance is highly dependent on the calibration set size, since the joint BatchNorm alignment requires a sufficient volume of statistics to correctly normalize the merged representation space.

\newpage
\subsection{Sensitivity to SRAC Temperature \texorpdfstring{$\beta$}{beta}}
Table~\ref{table-beta-sensitivity} summarizes the sensitivity of SRAC's soft routing scheme to the temperature hyperparameter $\beta$, as compared with our hyperparameter-free Minimalist Static Prototype Routing (MSPR).

\begin{table}[H]
\caption{Sensitivity of SRAC to temperature $\beta$ under WA on 256 calibration samples, compared to hyperparameter-free MSPR.}
\label{table-beta-sensitivity}
\vskip 0.05in
\begin{center}
\begin{footnotesize}
\setlength{\tabcolsep}{8pt}
\begin{tabular}{lcccc}
\toprule
Method ($\beta$) & MNIST & F-MNIST & CIFAR-10 & Average \\
\midrule
SRAC ($\beta=0.1$) & 52.86\% & 51.07\% & 23.46\% & 42.46\% \\
SRAC ($\beta=1.0$) & 55.96\% & 52.44\% & 24.76\% & 44.39\% \\
SRAC ($\beta=10.0$) & 70.71\% & 60.13\% & 32.75\% & 54.53\% \\
SRAC ($\beta=30.0$) & 74.09\% & 59.64\% & 39.15\% & 57.63\% \\
SRAC ($\beta=1000.0$) & 76.21\% & 56.07\% & 40.45\% & 57.58\% \\
\midrule
\textbf{MSPR (Ours, No $\beta$)} & \textbf{76.55\%} & \textbf{53.53\%} & \textbf{40.20\%} & \textbf{56.76\%} \\
\bottomrule
\end{tabular}
\end{footnotesize}
\end{center}
\vskip -0.15in
\end{table}

These results confirm that soft-routing methods are heavily reliant on meticulous hyperparameter tuning: with a suboptimal temperature, their performance degrades to levels significantly below our static, zero-tuning alternative.

\end{document}
